% =====================================================================================================
% Online Appendix A: Formal results and proofs (version 3, writer "appendices", 2026-09-24).
% v1: module m7_theory; v2: module ra5_theory; v3 (referee round 2): trimmed to the results the main text uses.
% Moved to the companion paper (paper/companion/proofs_from_appendixA_v2.tex): technical change over time (DMR),
% compute-equivalent gains, the Farrell decomposition, transmission, selection, allocation-based proxies, the
% growth-accounting corollaries, Lemma A3(v) and the frontier-scale illustration of the partial-identification result.
% New in v3: Proposition A4(vi) (binding data cap in a family; verified by code/paper/verify_family_cap.py); the
% on-path rate statement is a remark (Remark A2), not part of a proposition. Verification tallies and the claims
% registers: output/docs/verification.md. Notation (revision_plan_v3): allocation-error scale \varsigma; prices q_T, q_D,
% q_S; path slope e_p; 1-a written out (no b); transverse coordinate z; latency Z(N).
% Round 3 (WP3, fix list P1-P11): the value of compactness is the virtual value v_N = m_N + nu_N (Props. A-wedge and
% A-family); Corollary A-general (known costs that are not multiplicative); Prop. A-pi parts (vii)-(viii) (interval anchor,
% joint path bands, tilt allowance, union of anchored paths); A5 parts (iv) and (v) swapped to match Prop. 3; notation:
% upsilon (noise), varrho (price ratio), h_i and I (family), k in mathcal K, mathcal A (adoption), mathcal P and mathcal Q
% (Lemma A2), U_0 and V_0 (Prop. A1(iii)), theta_L and theta_U, e_L and e_U. Checks: code/analysis/rb5_theory/.
% =====================================================================================================

\section[Formal Results and Proofs]{Appendix \thesection. Formal Results and Proofs}\label{app:proofs}

\setcounter{lemma}{0}\setcounter{proposition}{0}\setcounter{corollary}{0}\setcounter{definition}{0}\setcounter{remark}{0}
\renewcommand{\thelemma}{A\arabic{lemma}}
\renewcommand{\theproposition}{A\arabic{proposition}}
\renewcommand{\thecorollary}{A\arabic{corollary}}
\renewcommand{\thedefinition}{A\arabic{definition}}
\renewcommand{\theremark}{A\arabic{remark}}
\renewcommand{\theHlemma}{A\arabic{lemma}}\renewcommand{\theHproposition}{A\arabic{proposition}}
\renewcommand{\theHcorollary}{A\arabic{corollary}}\renewcommand{\theHdefinition}{A\arabic{definition}}\renewcommand{\theHremark}{A\arabic{remark}}
\providecommand{\Var}{\operatorname{Var}}
\providecommand{\Cov}{\operatorname{Cov}}
\providecommand{\E}{\mathbb{E}}
\setlength{\emergencystretch}{2em}% allow looser lines in the online appendix

This appendix states and proves the formal results that the main text uses. All statements are exact unless an order of
approximation is given. Lemma~\ref{lemma:sigma} of the main text combines Lemmas~\ref{lem:sigma} and~\ref{lem:soc} with
Corollary~\ref{cor:complements}. The duality statements of Section~\ref{sec:framework} are Lemma~\ref{lem:alloc}.
Proposition~\ref{prop:wedge} is Proposition~\ref{prop:A-wedge}, with Proposition~\ref{prop:A-family} for family token
budgets and binding data caps, Lemmas~\ref{lem:geometry} and~\ref{lem:ce}, and Corollary~\ref{cor:suff}.
Proposition~\ref{prop:modelfree}(i)--(iv) is Proposition~\ref{prop:A-modelfree}(i)--(iv), and its part (v) is
Proposition~\ref{prop:A-pi}(ii)--(iv) and (vii)--(viii); Proposition~\ref{prop:ident}(i) combines
Proposition~\ref{prop:fd}(i), (ii) and (iv), part (ii) is Proposition~\ref{prop:fd}(iii), and part (iii) is
Proposition~\ref{prop:A-info}; the rate statement is Remark~\ref{rem:A-rate}. The statement of Section~\ref{sec:ident}
for costs that are not multiplicative is Corollary~\ref{cor:A-general}. Every algebraic identity below was also verified symbolically, and every
comparative-static, bias and information statement numerically, by brute-force optimization, simulation or direct
computation of Fisher information; the replication package lists each check and its tolerance
(\texttt{output/docs/verification.md}). Results on observational production functions and technical change over time
are in the companion paper \citep{okar2026observational}. Throughout, $E$ denotes irreducible loss and $\E$ the expectation operator.

\subsection{Setup and Notation}\label{app:setup}

\begin{definition}[Technology]\label{def:tech}
A training run uses $N$ parameters and $D$ training tokens (tokens processed, including repetitions) and attains a
held-out cross-entropy
\begin{equation}\label{eq:A-tech}
L \;=\; E + e^{-\omega}\Big[A\big(e^{\psi_N}N\big)^{-\alpha} + B\big(e^{\psi_D}D\big)^{-\beta}\Big]e^{\upsilon},
\qquad A,B,\alpha,\beta>0 .
\end{equation}
Lower-case letters are logs: $n=\ln N$, $d=\ln D$, $c=\ln C$ with training compute $C=6ND$, so $c=\ln 6+n+d$; $M\equiv D/N$.
Hicks-neutral productivity $\omega$ and factor-augmenting productivity $(\psi_N,\psi_D)$ are known to the lab when it
chooses inputs; the seed and evaluation noise $\upsilon$ is not. Write $u\equiv Ae^{-\alpha(n+\psi_N)}$, $v\equiv Be^{-\beta(d+\psi_D)}$ and
$R\equiv u+v$, so that $L-E=e^{-\omega+\upsilon}R$. The log output index is $y\equiv-\ln(L-E)$, and
$F(n,d)\equiv-\ln(Ae^{-\alpha n}+Be^{-\beta d})$. The output elasticities are
$\varepsilon_N\equiv-\partial\ln(L-E)/\partial n$ and $\varepsilon_D\equiv-\partial\ln(L-E)/\partial d$, and the wedge is
$w\equiv\varepsilon_N/\varepsilon_D$. Expansion-path objects are $a\equiv\beta/(\alpha+\beta)$ (so that
$1-a=\alpha/(\alpha+\beta)$), $\gamma\equiv\alpha\beta/(\alpha+\beta)$, $G\equiv(\alpha A/\beta B)^{1/(\alpha+\beta)}$,
$K\equiv\frac{\alpha+\beta}{\beta}AG^{-\alpha}$ and $\sigma^*\equiv2/(2+\alpha+\beta)$. Finally
$\chi\equiv\beta\psi_D-\alpha\psi_N$ is the bias index of productivity and $\Omega\equiv\omega+\gamma(\psi_N+\psi_D)$ is
productivity in compute-equivalent units.
\end{definition}

\begin{remark}[Exact log form]\label{rem:exact}
Taking logs of \eqref{eq:A-tech} gives, exactly and for any known $E$, $y=\omega+F(n+\psi_N,d+\psi_D)-\upsilon$,
$\varepsilon_N=\alpha u/(u+v)$ and $\varepsilon_D=\beta v/(u+v)$. None of the three depends on $\omega$, $\upsilon$ or
$E$; the elasticities use the $\omega$-free sum $u+v$, not the observed $L-E$.
\end{remark}

\begin{definition}[Generalized family]\label{def:kappa}
The generalized family replaces the bracket in \eqref{eq:A-tech} by
$\big[Ae^{-a_1(n+\psi_N)}+Be^{-b_1(d+\psi_D)}\big]^{\kappa}$ with $a_1,b_1,\kappa>0$. Chinchilla is $\kappa=1$
($a_1=\alpha$, $b_1=\beta$). The joint law of \citet[eq.~1.5]{kaplan2020scaling}, fitted to early-stopped runs with $D$ unique tokens and
non-embedding $N$, is the member with $E=0$, $a_1=\alpha_N/\alpha_D\approx0.738$, $b_1=1$ and $\kappa=\alpha_D\approx0.103$. For this family write $Q\equiv u+v$
with $u=Ae^{-a_1n}$ and $v=Be^{-b_1d}$, $S\equiv a_1+b_1$, $a\equiv b_1/S$ (so $1-a=a_1/S$), $\gamma\equiv\kappa a_1b_1/S$,
$G\equiv(a_1A/b_1B)^{1/S}$, $K\equiv\big(\tfrac{S}{b_1}AG^{-a_1}\big)^{\kappa}$ and $\sigma^*\equiv2/(2+S)$; with $\kappa=1$
these coincide with Definition~\ref{def:tech}. Because $R=Q^\kappa$ and the path does not depend on $\kappa$, the outer exponent leaves the expansion path
unchanged, multiplies the inner frontier elasticity $a_1b_1/S$ by $\kappa$ and raises the inner frontier constant to
the power $\kappa$ (Lemma~\ref{lem:alloc}(ii)).
\end{definition}

\subsection{The Elasticity of Substitution and the Second-Order Condition}\label{app:sigma}

\begin{lemma}[Elasticity of substitution; Lemma~\ref{lemma:sigma}]\label{lem:sigma}
For the technology \eqref{eq:A-tech}, and for the generalized family with $(a_1,b_1)$ in place of $(\alpha,\beta)$:
\begin{enumerate}
\item[(i)] The Hicks elasticity of substitution between $N$ and $D$ is
$\sigma=(\alpha u+\beta v)/[\alpha u(1+\beta)+\beta v(1+\alpha)]$.
\item[(ii)] $\sigma$ depends on $(N,D)$ only through the wedge,
\[
\sigma(w)=\frac{1+w}{1+\alpha+w(1+\beta)},\qquad\text{equivalently}\qquad\frac1\sigma=1+\frac{\alpha+\beta w}{1+w},
\]
which is monotone in $w$ between $1/(1+\alpha)$ ($w\to0$) and $1/(1+\beta)$ ($w\to\infty$).
\item[(iii)] On the compute-optimal expansion path ($w=1$), $\sigma=\sigma^*=2/(2+\alpha+\beta)$.
\item[(iv)] $\sigma$ is invariant to any strictly increasing transformation of output; in particular it does not depend
on $E$, $\omega$, $\upsilon$ or $\kappa$, or on whether output is measured by loss, bits per byte or a monotone
benchmark link.
\item[(v)] For the \citet{kaplan2020scaling} law, $\sigma\in[1/2,\,1/1.738]=[0.500,0.575]$ and $\sigma^*=2/3.738=0.535$.
\end{enumerate}
With two inputs the Hicks, Allen--Uzawa and Morishima elasticities coincide \citep{blackorby1989real}.
\end{lemma}

\begin{proof}
(i) With output $f(N,D)=-(AN^{-\alpha}+BD^{-\beta})$ in levels, $f_N=\alpha u/N$, $f_D=\beta v/D$,
$f_{NN}=-\alpha(1+\alpha)u/N^2$, $f_{DD}=-\beta(1+\beta)v/D^2$ and $f_{ND}=0$. The Hicks elasticity for two inputs,
$\sigma=-f_Nf_D(Nf_N+Df_D)/\{ND(f_{NN}f_D^2-2f_{ND}f_Nf_D+f_{DD}f_N^2)\}$, then equals the stated expression.
(ii) Substitute $\alpha u=w\beta v$; the derivative of $\sigma(w)$ has the sign of $\alpha-\beta$. (iii) is (ii) at
$w=1$. (iv) For $g$ strictly increasing, the first derivatives of $g\circ f$ are $g'$ times those of $f$, and the second
derivatives add terms proportional to $g''f_if_j$; in the denominator these sum to
$g''g'^2(f_N^2f_D^2-2f_N^2f_D^2+f_D^2f_N^2)=0$, so numerator and denominator both scale by $g'^3$. (v) Apply (ii) with
$(a_1,b_1)=(0.738,1)$.
\end{proof}

\begin{lemma}[Log coordinates and the second-order condition; Lemma~\ref{lemma:sigma}]\label{lem:soc}
Let $y=F(n,d)$ be twice continuously differentiable with $F_n,F_d>0$, and define $\mathcal P\equiv F_nF_d(F_n+F_d)>0$
and $\mathcal Q\equiv-(F_{nn}F_d^2-2F_{nd}F_nF_d+F_{dd}F_n^2)$. Then
\begin{enumerate}
\item[(i)] $\sigma=\mathcal P/(\mathcal P+\mathcal Q)$; hence $\sigma\in(0,1)$ if and only if $\mathcal Q>0$, that is, if and only if the isoquant is
strictly convex in $(\ln N,\ln D)$ coordinates.
\item[(ii)] Consider $\max y$ subject to $n+d=c-\ln6$ (equivalently $\min 6ND$ subject to $y\ge\bar y$). At a point with
$F_n=F_d$, the second derivative of $y$ along the isocost is $F_{nn}-2F_{nd}+F_{dd}=-\mathcal Q/F_n^2$. The point is a strict
local maximum of output (minimum of compute) if $\sigma\in(0,1)$, and a strict local minimum of output if $\mathcal Q<0$, which
requires $\sigma>1$ or $\sigma<0$.
\end{enumerate}
\end{lemma}

\begin{proof}
(i) With $f(N,D)=F(\ln N,\ln D)$: $f_N=F_n/N$, $f_D=F_d/D$, $f_{NN}=(F_{nn}-F_n)/N^2$, $f_{DD}=(F_{dd}-F_d)/D^2$ and
$f_{ND}=F_{nd}/(ND)$. Substituting into the Hicks formula, the powers of $N$ and $D$ cancel and
$\sigma=-F_nF_d(F_n+F_d)/[(F_{nn}-F_n)F_d^2-2F_{nd}F_nF_d+(F_{dd}-F_d)F_n^2]=\mathcal P/(\mathcal P+\mathcal Q)$. The
curvature of the log-isoquant $d=\phi(n)$ is $\phi''=\mathcal Q/F_d^3$, so $\mathcal Q>0$ is strict convexity. (ii) Along $(n,d)=(n_0-t,d_0+t)$,
$d^2y/dt^2=F_{nn}-2F_{nd}+F_{dd}$; at $F_n=F_d$ the bracket in $\mathcal Q$ equals $F_n^2(F_{nn}-2F_{nd}+F_{dd})$. If $\mathcal Q<0$,
then $\sigma=\mathcal P/(\mathcal P+\mathcal Q)$ lies outside $[0,1]$.
\end{proof}

\begin{corollary}[Complementarity, $\sigma<1$; Lemma~\ref{lemma:sigma}]\label{cor:complements}
Under the multiplicative cost $C=6ND$, every interior compute-optimal allocation that satisfies the second-order
condition strictly has $\sigma<1$ at the optimum. If $\sigma>1$ everywhere, compute-optimal training puts all compute on
one input, subject to bounds: for the CES technology $Y=(N^r+D^r)^{1/r}$ with $r\in(0,1)$, output on the isocost $ND=X$ is
$\sqrt X\,(2\cosh(rt))^{1/r}$ with $t=\tfrac12\ln(D/N)$, which is minimized, not maximized, at the symmetric point. Every
member of the generalized family has $\sigma<1$ everywhere.
\end{corollary}

\begin{proof}
The first two statements are Lemma~\ref{lem:soc}(ii); the CES computation is direct. For the generalized family,
Lemma~\ref{lem:sigma}(ii) gives $\sigma\le1/(1+\min\{a_1,b_1\})<1$.
\end{proof}

\begin{remark}
With a linear cost the second-order condition requires only $\sigma>0$; with the multiplicative cost the isocost is a line
in logs, so the isoquant must be convex in logs, $\sigma<1$; for any known cost the condition is $\sigma<\sigma_C$, the
elasticity of substitution of the isocost itself (Corollary~\ref{cor:A-general}). The empirical content of Corollary~\ref{cor:complements} is
that IsoFLOP profiles are U-shaped in $\ln N$; Proposition~\ref{prop:A-modelfree} turns the curvature of the U into the
level of $\sigma^*$.
\end{remark}

\subsection{Duality and Allocation}\label{app:duality}

\begin{lemma}[Cost function and conditional factor demands; the duality statements of Section~\ref{sec:framework}]\label{lem:alloc}
Consider the problem of minimizing $6ND$ subject to $L\le\bar L$, or equivalently of maximizing $y$ at given $c$, for
the technology \eqref{eq:A-tech}.
\begin{enumerate}
\item[(i)] The first-order condition is $\alpha u=\beta v$ ($\varepsilon_N=\varepsilon_D$), and it is sufficient because
$\sigma<1$ everywhere.
\item[(ii)] With $\psi_N=\psi_D=0$: $N^*=G(C/6)^{a}$, $D^*=G^{-1}(C/6)^{1-a}$ and $R^*(C)=K(C/6)^{-\gamma}$ with
$K=\frac{\alpha+\beta}{\beta}AG^{-\alpha}=(\alpha+\beta)(A/\beta)^{a}(B/\alpha)^{1-a}$. The cost function is
$C^*(R)=6(K/R)^{1/\gamma}$, with $\partial\ln C^*/\partial\ln R=-1/\gamma$; along the path
$\varepsilon_N=\varepsilon_D=\gamma$. For the generalized family the same statements hold with $(a_1,b_1)$ in place of
$(\alpha,\beta)$ and with $\gamma$ and $K$ of Definition~\ref{def:kappa}: $R^*(C)=K(C/6)^{-\gamma}$.
\item[(iii)] In general,
\begin{align*}
n^*&=\ln G+a(c-\ln6)+\chi/(\alpha+\beta),\quad y^*=\Omega+\gamma(c-\ln6)-\ln K,\\
d^*&=-\ln G+(1-a)(c-\ln6)-\chi/(\alpha+\beta),\\
\ln M^*(C)&=-2\ln G+(1-2a)(c-\ln6)-2\chi/(\alpha+\beta).
\end{align*}
\item[(iv)] At given compute, $\omega$ does not enter $(n^*,d^*)$; $\partial n^*/\partial\psi_D=a>0$,
$\partial d^*/\partial\psi_D=-a<0$ and $\partial\ln M^*/\partial\psi_D=-2a<0$: better data raise parameters and lower
both tokens and tokens per parameter. Only the bias index $\chi$ moves the mix; productivity with $\chi=0$ is equivalent
to a Hicks-neutral shift $\Delta\omega=\alpha\psi_N$.
\end{enumerate}
\end{lemma}

\begin{proof}
(i) The Lagrangian of $\max\{-\ln(u+v)\}$ subject to $n+d=c-\ln6$ gives $\alpha u=\beta v$. The upper contour sets of
$y$ are convex in $(n,d)$ (log-sum-exp is convex) and the constraint is linear. (ii) Substituting
$n=\ln G+a\ln(C/6)$ and $d=-\ln G+(1-a)\ln(C/6)$ gives $\alpha u/(\beta v)=(\alpha A/\beta B)G^{-(\alpha+\beta)}
(C/6)^{-\alpha a+\beta(1-a)}=1$, because $\alpha a=\beta(1-a)=\gamma$. At the optimum $v=\alpha u/\beta$, so
$R^*=u(\alpha+\beta)/\beta$; inverting $R^*(C)$ gives $C^*(R)$, and on the path $\varepsilon_N=\alpha u/R^*=\gamma$. For the
generalized family, apply this to the inner aggregate $Q$, whose path does not involve $\kappa$, and raise its frontier to
the power $\kappa$.
(iii) Effective inputs $n+\psi_N$ and $d+\psi_D$ enter $R$ exactly as $(n,d)$ do when $\psi=0$ and satisfy
$(n+\psi_N)+(d+\psi_D)=c-\ln6+\psi_N+\psi_D$. Apply (ii) with effective compute, subtract $\psi_N$ and $\psi_D$, and use
$a\psi_D-(1-a)\psi_N=\chi/(\alpha+\beta)$. (iv) Differentiate (iii). If $\alpha\psi_N=\beta\psi_D$ then
$u+v=e^{-\alpha\psi_N}(Ae^{-\alpha n}+Be^{-\beta d})$.
\end{proof}

Part~(ii) links the frontier $R^*(C)$ (a cost function), the demands $N^*(C),D^*(C)$ and the primal $L(N,D)$:
\citet{hoffmann2022training}'s Approaches~1 and~2 estimate the first two, and Approach~3 the primal.

\subsection{Geometry and the Wedge as a Sufficient Statistic}\label{app:geometry}

\begin{lemma}[Geometry of the technology; used in Propositions~\ref{prop:wedge}, \ref{prop:ident} and~\ref{prop:modelfree}]\label{lem:geometry}
For the generalized family (in effective inputs):
\begin{enumerate}
\item[(i)] $\nabla^2_{(n,d)}\ln R=\kappa\,\dfrac{uv}{Q^2}\begin{pmatrix}a_1^2&-a_1b_1\\-a_1b_1&b_1^2\end{pmatrix}$ has rank
one, and its null vector $(b_1,a_1)$ is the direction of the expansion path,
$(\partial n^*/\partial c,\partial d^*/\partial c)=(a,1-a)$.
\item[(ii)] With the transverse coordinate $z\equiv a_1n-b_1d$,
$\ln R=\kappa\big[-\tfrac12(a_1n+b_1d)+g(z)\big]$ with $g(z)=\ln(Ae^{-z/2}+Be^{z/2})$ and $g''(z)=s_u(1-s_u)\in(0,\tfrac14]$,
$s_u=u/Q$. Thus $\ln R$ is affine along every line with direction $(b_1,a_1)$, and $w$ and $\sigma$ are constant along
such lines.
\item[(iii)] $\ln w=z^*-z$ with $z^*=\ln(a_1A/b_1B)=S\ln G$; the expansion path is the level set $\{z=z^*\}$; and for any
run with compute $C$ and ratio $M$,
\begin{equation}\label{eq:A-wM}
\ln w=\tfrac{S}{2}\,\ln\!\big(M/M^*(C)\big)=\Big(\frac1{\sigma^*}-1\Big)\ln\!\big(M/M^*(C)\big),
\end{equation}
where $M^*(C)=G^{-2}(C/6)^{1-2a}$ is the training-optimal ratio at the run's own compute.
\item[(iv)] (Quasi-homotheticity.) $R(\lambda^{b_1}N,\lambda^{a_1}D)=\lambda^{-\kappa a_1b_1}R(N,D)$ for every $\lambda>0$.
More generally, for any $C^2$ technology, translation along $(v_n,v_d)$ maps isoquants onto isoquants if and only if
$\ln w$ is constant along it; then, with a unique optimum at each compute, the expansion path is a straight line in logs
with slope $a=v_n/(v_n+v_d)$, and all information about $\sigma$ is transverse to the path.
\end{enumerate}
\end{lemma}

\begin{proof}
(i) $\partial\ln Q/\partial n=-a_1u/Q$, so $\partial^2\ln Q/\partial n^2=a_1^2uv/Q^2$; the other entries follow in the same
way, and $\ln R=\kappa\ln Q$. The determinant is zero and $(b_1,a_1)$ is annihilated. (ii) Since
$-\frac12(a_1n+b_1d)\mp\frac12z$ equals $-a_1n$ and $-b_1d$ respectively, $Q=e^{-(a_1n+b_1d)/2}(Ae^{-z/2}+Be^{z/2})$; along
$(n_0+b_1t,d_0+a_1t)$, $z$ is constant and $a_1n+b_1d$ is affine in $t$. (iii) $w=a_1u/(b_1v)=(a_1A/b_1B)e^{-z}$. At given
compute write $n=n^*-x/2$ and $d=d^*+x/2$ with $x=\ln(M/M^*(C))$; then $z-z^*=-Sx/2$. (iv) The first statement is direct.
The log-isoquant has slope $-F_n/F_d$, so translates of isoquants are isoquants if and only if $F_n/F_d$ is constant
along the translation; with $w=1$ on the whole line through an optimum, which crosses every isocost once, that line
contains the optimum at every compute.
\end{proof}

\begin{lemma}[Allocative loss is a function of the wedge; used in Propositions~\ref{prop:wedge} and~\ref{prop:ident}]\label{lem:ce}
Let $C_{\min}(R)$ denote the least training compute $6ND$ that attains reducible loss $R$. For any run with wedge $w$,
\begin{equation}\label{eq:A-ce}
\frac{C}{C_{\min}(R)}=\Big(\frac{a_1+b_1w}{a_1+b_1}\Big)^{S/(a_1b_1)}w^{-1/a_1}
\;\;\overset{\kappa=1}{=}\;\;\Big(\frac{\alpha+\beta w}{\alpha+\beta}\Big)^{1/\gamma}w^{-1/\alpha},
\end{equation}
which depends neither on $(A,B,E,\kappa)$ nor on productivity. To second order,
$\ln(C/C_{\min})=(\ln w)^2/(2S)+O(|\ln w|^3)=\sigma^*(\ln w)^2/[4(1-\sigma^*)]+O(|\ln w|^3)$: the allocative loss is a
Harberger triangle whose area is governed by the elasticity of substitution.
\end{lemma}

\begin{proof}
Let $Q_0=u+v$ at the run. The cheapest bundle producing $Q_0$ has $a_1u^*=b_1v^*$, so $u^*=aQ_0$ and $v^*=(1-a)Q_0$.
Because $u\propto N^{-a_1}$ and $v\propto D^{-b_1}$, $C/C_{\min}=(u^*/u)^{1/a_1}(v^*/v)^{1/b_1}$ with $u/Q_0=b_1w/(a_1+b_1w)$,
which gives \eqref{eq:A-ce}. For the expansion, let $x=\ln w$: $\ln\big((a_1+b_1e^x)/S\big)=ax+\tfrac12a(1-a)x^2+O(x^3)$;
multiplying by $S/(a_1b_1)=1/[a(1-a)S]$ and subtracting $x/a_1=x/[(1-a)S]$ leaves $x^2/(2S)$. Finally
$S=2(1-\sigma^*)/\sigma^*$.
\end{proof}

\subsection{What On-Path Data Identify}\label{app:onpath}

\begin{definition}[On-path sample]\label{def:onpath}
A sample of runs is on-path if (a) all runs share the technology (generalized family, common parameters); (b)
productivity is Hicks-neutral ($\psi\equiv0$); (c) every run minimizes training compute for its loss (in the notation of
Proposition~\ref{prop:A-wedge}, $\delta=m_D=v_N=0$ and no constraint binds, so that $w=1$); and (d) compute varies over an interval (or at least five distinct values).
\end{definition}

\begin{proposition}[What on-path data identify; Proposition~\ref{prop:ident}]\label{prop:fd}
In an on-path sample:
\begin{enumerate}
\item[(i)] $n_i=\ln G+a(c_i-\ln6)$ and $d_i=-\ln G+(1-a)(c_i-\ln6)$. The inputs are deterministic affine functions of
compute, and the choice data reveal $a$ and $G$ (the observed compute-optimal ratio $M^*(C)$).
\item[(ii)] $\E[y\mid c]=\gamma(c-\ln6)-\ln K+\E[\omega\mid c]$. Without instruments for compute or panel restrictions
on $\omega$, the frontier elasticity $\gamma$ is identified if and only if $\E[\omega\mid c]$ is known up to an additive
constant, as in a single-lab sweep.
\item[(iii)] Suppose in addition $\E[\omega\mid c]$ is constant and $\kappa=1$. Write $c'\equiv c-\ln6$, $r_N\equiv\alpha a$
and $r_D\equiv\beta(1-a)$ for the rates at which the two terms decay along the observed path, and $U_0\equiv Ae^{-\alpha n_0}$,
$V_0\equiv Be^{-\beta d_0}$ (with $K$ normalized by $\E[e^{-\omega+\upsilon}]=1$), so that on the path
$\E[L\mid c]=E+U_0e^{-r_Nc'}+V_0e^{-r_Dc'}$, a two-rate exponential mixture whose rates coincide at the truth.
\begin{enumerate}
\item[(a)] (Global identification, by functional form.) $E$, $\gamma$ and $K$ are identified, and so are $\alpha=\gamma/a$
and $\beta=\gamma/(1-a)$: every parameter with $A,B>0$ that generates the same on-path outcomes has $r_N=r_D=\gamma$. The
ratio $A/B$ is identified only by imposing that the observed path is optimal ($G^{\alpha+\beta}=\alpha A/\beta B$);
outcomes identify only $U_0+V_0$.
\item[(b)] (No local identification.) The Jacobian of the on-path mean of $\ln L$ with respect to $(E,\ln A,\ln B,\alpha,\beta)$
has rank three. Because $A/B$ is not identified, every perturbation that moves $r_N$ and $r_D$ in opposite directions
leaves the on-path mean unchanged to first order (the change is absorbed by re-splitting $U_0+V_0$; at a given split
only one such direction is flat), so along each such
direction the least-squares or likelihood criterion is of fourth order in the step, and $\sigma^*$ is identified only at
second order whether or not $\alpha=\beta$. The twin map
$(\alpha,\beta,A,B)\mapsto\big(\beta(1-a)/a,\ \alpha a/(1-a),\ Be^{-\beta d_0+\beta(1-a)n_0/a},\ Ae^{-\alpha n_0+\alpha ad_0/(1-a)}\big)$
leaves the on-path outcome distribution exactly unchanged, so every parameter value off the ray $\alpha a=\beta(1-a)$ has
an observationally equivalent twin whose $\sigma^*$ differs unless $a=\tfrac12$. On the boundary of the parameter space
identification fails altogether: as $A\to0$ ($B\to0$) the on-path distribution converges to the true one for any
$\alpha$ ($\beta$).
\end{enumerate}
\item[(iv)] Without the restriction $\kappa=1$: for every $S>0$, the member
\begin{equation}\label{eq:A-family}
\begin{gathered}
a_1=(1-a)S,\qquad b_1=aS,\qquad \kappa_S=\frac{\gamma}{a(1-a)S},\\
A_S=\frac{b_1}{S}K^{1/\kappa_S}G^{a_1},\qquad B_S=\frac{a_1A_S}{b_1G^{S}}
\end{gathered}
\end{equation}
generates exactly the same joint distribution of $(n,d,y)$ as the true technology. Hence $\sigma^*=2/(2+S)$ is not
identified: every value in $(0,1)$ is consistent with the data, whereas by Lemma~\ref{lem:soc} no value $\sigma^*>1$ is (in the
generalized family, none with $\sigma^*\ge1$).
\end{enumerate}
\end{proposition}

\begin{proof}
(i) is Lemma~\ref{lem:alloc}(ii). (ii) On the path $y=\omega+\gamma(c-\ln6)-\ln K-\upsilon$; if $\E[\omega\mid c]=h(c)$
with $h$ unknown, then $\gamma(c-\ln6)+h(c)$ and $\tilde\gamma(c-\ln6)+h(c)+(\gamma-\tilde\gamma)(c-\ln6)$ are
indistinguishable. (iii)(a) With $\omega$ independent of $c$, $\E[L\mid c]=E+\tilde Ke^{-\gamma c'}$; a three-parameter
exponential is identified from an interval of $c$. If $(\tilde\alpha,\tilde\beta,\tilde A,\tilde B)$ with $\tilde A,\tilde
B>0$ fits the same curve on the path $n=n_0+ac'$, $d=d_0+(1-a)c'$, then
$\tilde U_0e^{-\tilde\alpha ac'}+\tilde V_0e^{-\tilde\beta(1-a)c'}=\tilde Ke^{-\gamma c'}$ on an interval; exponentials with
distinct rates are linearly independent, so $\tilde\alpha a=\tilde\beta(1-a)=\gamma$. The remaining condition,
$\tilde U_0+\tilde V_0=\tilde K$, is one equation in two unknowns. (b) On the path $u=u_0e^{-\gamma c'}$ and
$v=v_0e^{-\gamma c'}$. The scores of $\ln L$ are $1/L$ times: $1$ for $E$; $u$ and $v$ for $\ln A$ and $\ln B$;
$-(n_0+ac')u$ for $\alpha$; and $-(d_0+(1-a)c')v$ for $\beta$. They span $\{1,e^{-\gamma c'},c'e^{-\gamma c'}\}$, so the rank is
three. Every $(U_0,V_0)$ with $U_0+V_0=U_0^\circ+V_0^\circ$ (the true values) and equal rates generates the true
distribution, and at such a point the first-order change of the mean from $(\delta r_N,\delta r_D)$ is
$-c'e^{-\gamma c'}(U_0\delta r_N+V_0\delta r_D)$, which vanishes for $U_0/V_0=-\delta r_D/\delta r_N>0$ whenever
$\delta r_N\delta r_D<0$; at a given split $U_0/V_0$ this holds for one direction only. The second-order term,
$\tfrac12c'^2e^{-\gamma c'}(U_0\delta r_N^2+V_0\delta r_D^2)$, is not in the span of $\{1,e^{-\gamma c'},c'e^{-\gamma c'}\}$, so
the criterion is exactly of fourth order along such directions. Since $\delta(\alpha+\beta)=\delta r_N/a+\delta r_D/(1-a)$,
generic opposite-sign directions move $\sigma^*$ at first order. The twin claim follows by substitution: the twin's two
terms are $V_0e^{-r_Dc'}$ and $U_0e^{-r_Nc'}$ with the roles of the inputs exchanged. (iv) Apply Lemma~\ref{lem:alloc}(ii) to the
inner aggregator of member $S$: its path slope is $b_1/S=a$, its path intercept is $G$ by construction of $B_S$, and its
inner frontier is $K^{1/\kappa_S}(C/6)^{-a_1b_1/S}$ by construction of $A_S$. Hence $R^*=K(C/6)^{-\kappa_Sa_1b_1/S}=K(C/6)^{-\gamma}$,
and inputs and outputs coincide run by run for every $(\omega,\upsilon)$.
\end{proof}

\begin{remark}[Rates and inference on the path]\label{rem:A-rate}
The on-path model of Proposition~\ref{prop:fd}(iii) is an over-fitted mixture: writing its mean as
$E+(U_0+V_0)\int e^{-rc'}\,d\mathcal G(r)$, with $\mathcal G$ the two-point distribution on $\{r_N,r_D\}$, the data
identify the mean of $\mathcal G$ through the coefficient on $c'e^{-\gamma c'}$ and its second moment about $\gamma$
through the coefficient on $c'^2e^{-\gamma c'}$, each at
best at rate $n^{-1/2}$, so the spread $|r_N-r_D|$, and with it $\sigma^*$, at best at rate $n^{-1/4}$, as in the bound
for over-fitted mixtures of \citet{chen1995optimal}. The bound concerns estimators that leave $\beta/(\alpha+\beta)$ free,
such as unrestricted least-squares fits of the Chinchilla form. Imposing $\beta/(\alpha+\beta)=a$, which the optimality
that defines an on-path sample implies under $\kappa=1$, forces $r_N=r_D$ and removes the singularity:
$\sigma^*=2/\{2+\gamma/[a(1-a)]\}$ combines $\hat a$ from the choices with $\hat\gamma$ from a regular three-parameter
exponential and is estimated at the usual root-$n$ rate. This is a heuristic argument, not a theorem for this model; the
results of \citet{sargan1983identification} and \citet{rotnitzky2000likelihood} on singular information do not apply
directly because $U_0/V_0$ is not identified at the truth.

A Monte Carlo supports it. It holds the design fixed, 40 on-path compute levels spanning four decades, and lets the
standard deviation of Gaussian noise on $L$ fall from $10^{-3}$ to $10^{-7}$ in five steps, with 100 replications at each;
for least squares this is equivalent to replicating the design $n$ times, with $n$ proportional to the inverse noise
variance. The unrestricted fit is least squares over $(E,U_0,V_0,r_N,r_D)$, equivalently over $(E,\ln A,\ln B,\alpha,\beta)$,
with $U_0$ and $V_0$ bounded below at 10 percent of the smaller true value, which excludes the boundary where an exponent is
not identified; the error recorded is the larger of the errors of the two exact twins of Proposition~\ref{prop:fd}(iii)(b).
The Hoffmann et al.\ exponents ($\alpha\ne\beta$) and equal exponents give the same picture. The 75th and 90th
percentiles of $|\hat\sigma^*-\sigma^*|$ scale as $n^{-0.24}$ and as $n^{-0.27}$ to $n^{-0.28}$, and as $n^{-0.50}$ when
the restriction $\beta/(\alpha+\beta)=a$ is imposed; the 25th percentile scales as $n^{-1/2}$, because in about 40 percent
of replications the fitted rates coincide. Wald and delta-method intervals based on the near-singular Hessian are then
invalid (Online Appendix~\ref{app:mc}).
\end{remark}

\begin{proposition}[Information rates; Proposition~\ref{prop:ident}]\label{prop:A-info}
Consider a sweep with $y_i=F(n_i,d_i;\phi)-\upsilon_i$, $\upsilon_i$ i.i.d.\ $N(0,V_\upsilon)$, $E$ known, and the generalized
family parameterized by $\phi=(S,a,\gamma,\ln G,\ln K)$. For run $i$ let $w_i$ be its wedge,
$\Delta_i\equiv(1-a)n_i-ad_i-\ln G=-(\ln w_i)/S$ its transverse deviation (Lemma~\ref{lem:geometry}), and
$\Phi_i\equiv\ln(C_i/C_{\min}(R_i))$ its allocative loss (Lemma~\ref{lem:ce}).
\begin{enumerate}
\item[(i)] (Exact decomposition.) $y_i=\gamma(c_i-\ln6)-\ln K-\gamma\,\Phi_S(\Delta_i)-\upsilon_i$, where
$\Phi_S(\Delta)=\big[\ln(1-a+ae^{-S\Delta})+aS\Delta\big]/[a(1-a)S]=\Phi_i\ge0$ and $\Phi_S(0)=\Phi_S'(0)=0$. The curvature
parameter $S$, and hence $\sigma^*$, enters output only through the runs' allocative losses.
\item[(ii)] (Scores.) $\partial y_i/\partial\ln G=\gamma(1-w_i)/(1-a+aw_i)=-\gamma\ln w_i+O\big((\ln w_i)^2\big)$ and
$\partial y_i/\partial S=-(\gamma/S)\,\Phi_i\,[1+O(|\ln w_i|)]=-\gamma(\ln w_i)^2/(2S^2)+O(|\ln w_i|^3)$.
\item[(iii)] (Efficient information.) Let $\zeta_k$ be the vector of scores for parameter $k$ and $M_{-k}$ the
residual-maker for the other four score vectors. Then $\mathcal I^*(\ln G)=(\gamma^2/V_\upsilon)\,\|M_{-G}\ln w\|^2\,
[1+O(\max_i|\ln w_i|)]$ and $\mathcal I^*(S)\le\big(\gamma^2/(S^2V_\upsilon)\big)\sum_i\Phi_i^2\,[1+O(\max_i|\ln w_i|)]
=O\big(\sum_i(\ln w_i)^4\big)$. If every run lies on the expansion path, both informations are zero. If the transverse
deviations of a fixed design are scaled by $\varsigma\to0$, then $\mathcal I^*(\ln G)=O(\varsigma^2)$ and
$\mathcal I^*(\sigma^*)=O(\varsigma^4)$, so the standard error of $\hat\sigma^*$ grows at least like $\varsigma^{-2}$ and
that of $\ln\hat G$ at least like $\varsigma^{-1}$. The same order holds for the path's location at any given compute
$C_0$, $\ln M^*(C_0)=-2\ln G+(1-2a)\ln(C_0/6)$, whose efficient information combines $\zeta_G$ with the score for $a$;
both are $O(\varsigma)$, but the level of the information differs from $\mathcal I^*(\ln G)/4$. Both rates are attained,
so that the standard errors grow exactly like $\varsigma^{-2}$ and $\varsigma^{-1}$, by designs that cross several compute
levels with symmetric transverse offsets, such as factorial designs.
\item[(iv)] If $E$ is estimated and the noise is Gaussian on $\ln L$, the score for $S$ is multiplied by
$R_i/L_i\in(0,1)$, so the bound in (iii) still holds.
\end{enumerate}
\end{proposition}

\begin{proof}
(i) Because the frontier is a power law, $R=K(C_{\min}(R)/6)^{-\gamma}$ for every $R$, so
$\ln R=\ln K-\gamma(c-\ln6)+\gamma\ln(C/C_{\min}(R))$. By Lemmas~\ref{lem:ce} and~\ref{lem:geometry}(iii),
$\ln(C/C_{\min})=\ln\big((a_1+b_1w)/S\big)/[a(1-a)S]-\ln w/a_1$ with $w=e^{-S\Delta}$, $a_1=(1-a)S$ and $b_1=aS$, which is
$\Phi_S(\Delta)$; it is nonnegative because $C\ge C_{\min}$, and $\Phi_S'(\Delta)=(1-e^{-S\Delta})/(1-a+ae^{-S\Delta})$
vanishes at zero. (ii) Holding $(n,d)$ fixed, $\Delta$ depends on $\ln G$ (with derivative $-1$) but not on $S$, so
$\partial y/\partial\ln G=\gamma\Phi_S'(\Delta)$. Write $\Phi_S=h(S\Delta)/[a(1-a)S]$ with $h(x)=\ln(1-a+ae^{-x})+ax$; then
$\partial\Phi_S/\partial S=[xh'(x)-h(x)]/[a(1-a)S^2]$, and $h(x)$ and $xh'(x)-h(x)$ both equal $\tfrac12a(1-a)x^2+O(x^3)$.
(iii) Under Gaussian errors the efficient information for one parameter is the Schur complement
$V_\upsilon^{-1}\zeta_k'M_{-k}\zeta_k$; the inequality uses $\|M_{-S}\zeta_S\|\le\|\zeta_S\|$ and (ii), and the orders follow
from $\zeta_G=O(\varsigma)$ and $\zeta_S=O(\varsigma^2)$. For a design that crosses several compute levels with symmetric
transverse offsets, $\Delta_i^2$ is not in the span of the leading-order nuisance scores
$\{1,\,c_i-\ln6-\Phi_i,\,\Delta_i,\,\Delta_i(c_i-\ln6)\}$ and $\Delta_i$ is not in the span of the others, which gives
attainment; in a nine-budget, five-offset factorial design the numerical log-log slopes are 4.00 and 2.00, also with $E$
estimated, and 2.00 for $\ln M^*(C_0)$ at a fixed $C_0$. (iv) $\partial\ln L/\partial S=-(R/L)\,\partial y/\partial S$.
\end{proof}

\subsection{What Over-Training Reveals: The Generalized Wedge}\label{app:A-wedge-sub}

\begin{proposition}[Generalized wedge; Proposition~\ref{prop:wedge}]\label{prop:A-wedge}
A developer chooses $(N,D)$ to maximize
\begin{equation}\label{eq:A-gw-obj}
\begin{gathered}
\Pi(N,D)=V\big(L(N,D),N\big)-X_T(N,D)-q_DD\\
\textup{subject to}\quad g_k(n,d)\le0,\quad k\in\mathcal K,
\end{gathered}
\end{equation}
where $V$ is differentiable with $\partial V/\partial L<0$; $X_T\equiv q_T(N)\,6ND$ is training expenditure at the
(shadow) price $q_T(N)$ of a training FLOP, whose size elasticity is $\delta\equiv d\ln q_T/d\ln N$; $q_D\ge0$ is a cost
per training token; $L$ is given by \eqref{eq:A-tech} or by any technology with $\varepsilon_N,\varepsilon_D>0$; and the
$g_k$ are differentiable constraints. Let $m_N\equiv-(\partial V/\partial\ln N)_L/X_T$ be the objective's own value of
compactness at fixed quality per unit of training expenditure, $m_D\equiv q_DD/X_T$, and $\xi_k\ge0$ the multipliers
(value per log point). Let
\[
\nu_N\equiv\sum_{k\in\mathcal K:\,g_{k,d}=0}\xi_kg_{k,n}/X_T,\qquad v_N\equiv m_N+\nu_N:
\]
$\nu_N$ is the shadow value of the constraints on size alone per unit of training expenditure (a cap has $g_{k,n}>0$ and a
minimum-size floor $g_{k,n}<0$, so floors enter with a negative sign), and $v_N$ is the value of compactness, a virtual
price in the sense of \citet{neary1980theory}: the value of $m_N$ at which a developer free of those constraints would
make the same choice. Without binding constraints on size, $v_N=m_N$.
\begin{enumerate}
\item[(i)] At any interior Karush--Kuhn--Tucker point,
\begin{equation}\label{eq:A-gw}
w=\frac{1+\delta+m_N+\sum_k\xi_kg_{k,n}/X_T}{1+m_D+\sum_k\xi_kg_{k,d}/X_T}.
\end{equation}
With no binding constraint that involves $d$, $w=(1+\delta+v_N)/(1+m_D)$; if also $\delta=m_D=0$, $w-1=v_N$, and the
compactness share $s\equiv(w-1)/w$ is $v_N/(1+v_N)$.
\item[(ii)] (Invariance and contamination.) $w$ is computed from $(N,D)$ and the technology alone; it is invariant to
$\omega$, $\upsilon$, $E$ and $\kappa$. If the developer's first-order conditions involve a technology $L'$ (its own
productivity, a law it believes, or its own output concept) while the econometrician uses $L$, then
$\hat w=w_{L'}\cdot w_L/w_{L'}$, where $w_{L'}$ is the wedge that \eqref{eq:A-gw} determines at $(N,D)$ under $L'$ and
$w_L$ the one that $L$ assigns to the same inputs. With factor-augmenting productivity $w_L/w_{L'}=e^{-\chi}$, so
$\hat w=w_{L'}e^{-\chi}$.
\item[(iii)] (Special cases.) Let $T$ be lifetime inference tokens in present value (discounted at the interest rate plus
an obsolescence hazard), $q_S(N)$ the price of a serving FLOP, $X_S\equiv q_S\,2NT$ serving expenditure and
$\varrho\equiv q_S/q_T$ the price ratio.
\begin{enumerate}
\item[(a)] (Serving cost borne at rate $\lambda$.) If $V=R(L)-\lambda X_S$ with $T=T(L,N)$, $q_D=0$ and no binding
constraint, then
$w=1+\delta+\lambda\varrho(\phi+\eta)\,T/(3D)$, where $\phi\equiv1+d\ln q_S/d\ln N$ and
$\eta\equiv(\partial\ln T/\partial\ln N)_L$; a direct effect of size on revenue at fixed quality adds
$-(\partial R/\partial\ln N)_L/X_T$. With $\delta=\eta=0$ and $\phi=1$: $w=1+\lambda\varrho T/(3D)$, $w-1=\lambda X_S/X_T$
and $s=\lambda X_S/(X_T+\lambda X_S)$; no conversion of FLOPs into dollars is needed. Compute proportional to $N$ that is
planned at training time, such as post-training, reinforcement-learning rollouts or synthetic-data generation, enters in
the same way. Counting post-training expenditure $X_P$ as a cost alongside $X_T$ (so that $V$ excludes it), with
$x_P\equiv X_P/X_T$ and $\delta=m_D=0$, gives $w=1+x_P+m_N$, where $m_N$ is net of post-training.
\item[(a$'$)] (Adoption.) If the developer bears no serving cost and $V=v(\mathcal A)$, where adoption $\mathcal A(L,P_u)$
falls with users' serving cost per token $P_u\propto N$, then $m_N=v'(\mathcal A)\mathcal A\,|\varepsilon_{\mathcal
A,P}|/X_T$: $w-1$ is the adoption value of compactness, not a token count.
\item[(b)] (Tier or memory cap.) If $N\le\bar N$ binds, $w=(1+\delta+v_N)/(1+m_D)$ with $\nu_N=\xi/X_T\ge0$; $w-1$
identifies $v_N$ and bounds $m_N$ from above (when $\delta=m_D=0$), but does not identify $m_N$: any
$w\ge(1+\delta+m_N)/(1+m_D)$ is rationalizable, including with $m_N=0$. A menu of sizes is this case at the chosen size,
with $\nu_N$ of either sign (negative when the menu binds as a floor).
\item[(c)] (Latency.) A binding service-level objective $Z(N)\le\bar Z$, with elasticity $\ell\equiv d\ln Z/d\ln N>0$,
is case (b) with $\nu_N=\ell\xi_Z/X_T$. A latency cost inside $V$ instead enters $m_N$.
\item[(d)] (Data cost and data cap.) $w=(1+\delta+v_N)/(1+m_D)$, so a data cost lowers $w$ for given $v_N$. A binding cap
$D\le\bar D$ adds $\mu\equiv\xi_D/X_T\ge0$ to the denominator, $w=(1+\delta+v_N)/(1+m_D+\mu)$, so with $\delta=m_D=0$,
$w-1$ bounds $v_N$ from below; $w$ can fall below one with $v_N=0$.
\item[(e)] (Logit distillation.) If distillation leaves the student's technology unchanged, teacher forward passes cost
$2N_TD$ FLOPs at the training price, a per-token cost with $m_D=N_T/(3N)$: with $\delta=0$ and $\phi=1$,
$w=\big(1+\lambda\varrho T/(3D)\big)/\big(1+N_T/(3N)\big)$. If it also raises the productivity of tokens ($\chi>0$), the
measured wedge falls further, by (ii).
\item[(f)] (Utilization varying with size.) If the prices of training and serving FLOPs have size elasticities $\delta$
and $\delta_S$ ($\phi=1+\delta_S$), then $w=1+\delta+\lambda(1+\delta_S)\varrho T/(3D)$, which equals
$(1+\delta)\big(1+\lambda\varrho T/(3D)\big)$ if and only if $\delta_S=\delta$.
\end{enumerate}
\item[(iv)] (Signs, on $m_N$.) Let $w_0\equiv(1+\delta+m_N)/(1+m_D)$. A single binding constraint raises $w$ above $w_0$
if and only if $g_n>w_0\,g_d$. Relative to the objective's own $m_N$, memory or tier caps and latency objectives raise
$w$ (they raise $v_N$); data caps and deadlines that bind on tokens lower it; deadlines proportional to compute pull it
toward one; minimum-size floors lower it. Hence $\hat T\equiv3D(\hat w-1)$ overstates $\lambda\varrho T$ under caps,
latency objectives, $\delta>0$ (with $\delta_S\ge0$), post-training compute and $\chi<0$, and understates it under data
costs, data caps, distillation, minimum-size floors, $\delta<0$ (with $\delta_S\le0$) and $\chi>0$. A wedge below one
cannot be generated by serving demand as long as serving expenditure rises with size at given quality ($\phi+\eta\ge0$).
\end{enumerate}
\end{proposition}

\begin{proof}
In $(n,d)$ the Lagrangian is $\Pi-\sum_k\xi_kg_k$. Since $\partial L/\partial n=-(L-E)\varepsilon_N$ and
$\partial L/\partial d=-(L-E)\varepsilon_D$, with $\Lambda\equiv-(\partial V/\partial L)(L-E)>0$,
$\partial X_T/\partial n=(1+\delta)X_T$ and $\partial X_T/\partial d=X_T$, the first-order conditions are
$\Lambda\varepsilon_N+(\partial V/\partial\ln N)_L-(1+\delta)X_T-\sum_k\xi_kg_{k,n}=0$ and
$\Lambda\varepsilon_D-X_T-q_DD-\sum_k\xi_kg_{k,d}=0$; dividing gives \eqref{eq:A-gw}. Constraints with $g_{k,d}=0$ enter
only the numerator, where they add $\nu_N$ to $m_N$. (ii) Remark~\ref{rem:exact}; in the generalized family
$w=a_1u/(b_1v)$ does not involve $\kappa$. If the first-order conditions involve $L'$, \eqref{eq:A-gw} determines
$w_{L'}(N,D)$, and the econometrician computes $w_L(N,D)=w_{L'}(N,D)\cdot w_L/w_{L'}$; with factor-augmenting
productivity $w_{L'}/w_L=e^{\beta\psi_D-\alpha\psi_N}=e^{\chi}$. (iii)(a) $-(\partial V/\partial\ln N)_L=\lambda
X_S(\phi+\eta)$ because $\ln X_S=\ln q_S+\ln2+\ln N+\ln T$; divide by $X_T$ and use $X_S/X_T=\varrho T/(3D)$.
Post-training compute has elasticity one in $N$ and zero in $D$. (a$'$) $(\partial V/\partial\ln N)_L=v'(\mathcal
A)\,\partial\mathcal A/\partial\ln P_u$. (b) $g=n-\bar n$, so $g_n=1$ and $g_d=0$, and $\nu_N=\xi/X_T$. A menu fixes $n$ at
the chosen size; with $\xi\equiv\partial\Pi/\partial n$ there, of either sign, the condition in $n$ holds with
$\nu_N=\xi/X_T$. (c) $g=\ln Z(n)-\ln\bar Z$, so $g_n=\ell$ and $g_d=0$. (d) The cap is $g=d-\bar d$. (e) The teacher's cost
$2N_TDq_T$ has elasticity zero in $N$ and one in $D$, so $m_D=2N_TD/(6ND)$; a rise in the productivity of tokens is
$\chi>0$ in (ii). (f) The elasticity of $X_T$ in $N$ is $1+\delta$ and that of $X_S$ is $1+\delta_S$. (iv) With
$A_0\equiv1+\delta+m_N$ and $B_0\equiv1+m_D$, $\partial_\xi[(A_0+\xi g_n/X_T)/(B_0+\xi g_d/X_T)]$ has the sign of
$g_nB_0-g_dA_0$ for every $\xi$. The first-order conditions are necessary, and the inversion uses nothing else; they are
also sufficient when $\Pi$ is concave in $(n,d)$, as in the lifetime-compute problem. Every case was checked by brute-force
maximization of the objective, including binding size caps, floors, menus and latency objectives with and without a cap on
tokens (\texttt{code/analysis/rb5\_theory/verify\_virtual\_value.py}).
\end{proof}

\begin{remark}[Conduct and the parallel with markup estimation]\label{rem:A-conduct}
Cases (a), (a$'$) and (b) are three models of conduct: a developer that bears the cost of serving its own model; an
open-weight developer that internalizes users' costs only through adoption; and size-class competition with $N$ set by a
menu of deployment tiers \citep[compare the Leontief model of][in which over-training does not arise]{hao2026theory}.
Under (a), $w-1$ rises with the developer's own serving footprint; under (a$'$), with the sensitivity of adoption to users'
serving cost; under (b), sizes bunch at tier boundaries and $w-1$ is unrelated to demand given the tier. The wedge is the
ratio statistic of \citet{raval2023testing} with FLOP-accounting weights, and part (ii) contains the two classic critiques
of the production approach to markups, factor-augmenting heterogeneity \citep{demirer2020production} and a mismatch between
measured and valued output \citep{bond2020unpleasant}. Here the elasticities come from designed experiments rather than
from proxy estimation on the same choices \citep{doraszelski2021reexamining}. As in \citet{deridder2026hitchhiker}, levels
are fragile and ranks robust: at given compute the rank of $\hat w$ across models is the rank of $M$ for every member of the
generalized family.
\end{remark}

\begin{corollary}[Sufficient statistics; used in Proposition~\ref{prop:wedge}]\label{cor:suff}
In the generalized family, for a model with compute $C$ and ratio $M$ at an interior optimum of
Proposition~\ref{prop:A-wedge} with $\delta=m_D=0$ and no binding constraint on tokens,
\[
v_N=w-1=\Big(\frac{M}{M^*(C)}\Big)^{1/\sigma^*-1}-1,\qquad s=1-\Big(\frac{M}{M^*(C)}\Big)^{-(1/\sigma^*-1)},
\]
and $C/C_{\min}$ is given by Lemma~\ref{lem:ce}. $M^*(C)$ is the training-optimal ratio at the model's own compute. Under
Proposition~\ref{prop:A-wedge}(iii)(a) with $\phi=1$, $\eta=0$ and no binding size constraint,
$\lambda\varrho T/D=3[(M/M^*(C))^{1/\sigma^*-1}-1]$. In the homothetic case $\alpha=\beta=\rho$, $M^*$ is constant and
$C/C_{\min}=\cosh(\tfrac\rho2\ln(M/M^*))^{2/\rho}$. With unequal exponents, $\ln w=\tfrac{\alpha+\beta}{2}\ln(M/M^*(C))$
still holds, with the mean exponent $(\alpha+\beta)/2=1/\sigma^*-1$ in place of $\rho$ and $M^*(C)$ in place of $M^*$, but
$C/C_{\min}$ is then given by Lemma~\ref{lem:ce}: the $\cosh$ form holds only to second order in $\ln w$ (with the
exponents of \citet{hoffmann2022training}, it gives 78.0 against an exact 101.2 at $w=13.5$). This justifies evaluating a
technology's compute-optimal law at each model's own compute.
\end{corollary}

\begin{proof}
Combine Proposition~\ref{prop:A-wedge}(i), equation~\eqref{eq:A-wM} and Lemma~\ref{lem:ce}. Both
$\ln(C/C_{\min})$ from Lemma~\ref{lem:ce} and $\tfrac{2}{\rho}\ln\cosh(\tfrac\rho2\ln(M/M^*(C)))$ with
$\rho=(\alpha+\beta)/2$ equal $(\ln w)^2/[2(\alpha+\beta)]+O(|\ln w|^3)$.
\end{proof}

\begin{proposition}[Family token budgets and binding data caps; Proposition~\ref{prop:wedge}]\label{prop:A-family}
A developer trains $I$ models of sizes $N_1,\dots,N_I$ on a common token budget $D$ and chooses $D$ (and, where stated,
the sizes) to maximize $\sum_iV_i(L(N_i,D),N_i)-\sum_iX_{T,i}-q_DD$ subject to $\underline n_i\le n_i\le\bar n_i$ and
$d\le\bar d$, where $X_{T,i}=q_T\,6N_iD$ and the data cost is incurred once. Let $w_i$ be member $i$'s wedge computed from
the technology, $h_i\equiv N_i/\sum_jN_j$ its compute share, $s_i\equiv(w_i-1)/w_i$,
$m_{N,i}\equiv-(\partial V_i/\partial\ln N_i)_L/X_{T,i}$, $m_D\equiv q_DD/\sum_jX_{T,j}$,
$\Lambda_i\equiv-(\partial V_i/\partial L)(L_i-E)$, and $\rho_i\equiv\Lambda_i\varepsilon_{D,i}/X_{T,i}$, the marginal
value of data to member $i$ per unit of its marginal training cost. Let $\xi_i$ be the net multiplier on member $i$'s size
constraints (zero if they are slack, $\xi_i\ge0$ at a cap, $\xi_i\le0$ at a floor) and $v_{N,i}\equiv m_{N,i}+\xi_i/X_{T,i}$
its value of compactness (Proposition~\ref{prop:A-wedge}). Let $\bar w_H\equiv(\sum_ih_i/w_i)^{-1}$ be the
compute-weighted harmonic mean of the member wedges and $\pi_i\equiv(h_i/w_i)/\sum_j(h_j/w_j)$.
\begin{enumerate}
\item[(i)] (Member wedges are technological.) Given $D$, $\ln w_i-\ln w_j=\ln w(n_i,d)-\ln w(n_j,d)$, which equals
$-a_1(n_i-n_j)$ in the generalized family: the profile of $\hat T_i\equiv3D(w_i-1)$ across sizes is fixed by the technology.
\item[(ii)] (Family first-order condition.) At an interior optimum in $D$, $\sum_ih_i\rho_i=1+m_D$. If the sizes are
set by a menu, this is the only first-order condition: the $w_i$ carry no information about the $m_{N,i}$, and (iii)
applies to the virtual values $v_{N,i}$.
\item[(iii)] (Sizes chosen or constrained.) If each $n_i$ is at an interior optimum or at a binding size constraint,
$\sum_i\pi_i(1+v_{N,i})=(1+m_D)\,\bar w_H$ with $v_{N,i}\equiv m_{N,i}+\xi_i/X_{T,i}$, where $\bar w_H$ is computed from
the technology alone and satisfies $1-1/\bar w_H=\sum_ih_is_i$: the family's first-order condition pins the
$\pi$-weighted average of the members' values of compactness, whether the sizes are chosen at the margin or set by
binding size constraints. Under Proposition~\ref{prop:A-wedge}(iii)(a) with common $(\lambda,\varrho)$, $\delta=\eta=0$,
$\phi=1$, $m_D=0$ and no binding size constraint: $\sum_i\pi_i\hat T_i=\lambda\varrho\sum_i\pi_iT_i$. Only this average is
a revealed-preference object; the member values are not identified (one equation for $I$ unknowns).
\item[(iv)] (Binding size caps and floors.) Binding size caps make $\xi_i\ge0$ and binding minimum-size floors make
$\xi_i\le0$. If every size is at a binding cap, $\sum_i\pi_i(1+m_{N,i})\le(1+m_D)\bar w_H$: $\bar w_H-1$ bounds the
$\pi$-weighted average of the objective's own $m_{N,i}$ from above (with $m_D=0$). Binding floors reverse the inequality.
\item[(v)] (A common $D$ under member-level choice is a knife-edge.) If instead each member chooses its own $D_i$ with
$\delta=m_D=0$ and no binding constraint, equal token counts $D_i=\bar D$ occur only if $1+m_{N,i}=w(n_i,\bar d)$ for every
$i$; in the generalized family, only if $(1+m_{N,i})/(1+m_{N,j})=(N_i/N_j)^{-a_1}$ exactly.
\item[(vi)] (Binding data cap.) If the token budget is at a binding cap $d=\bar d$ with multiplier $\xi_D\ge0$ and each
$n_i$ is at an interior optimum or at a binding size constraint, then
$\sum_i\pi_i(1+v_{N,i})=(1+m_D+\mu)\,\bar w_H\ge(1+m_D)\,\bar w_H$ with $\mu\equiv\xi_D/\sum_jX_{T,j}$: with $m_D=0$,
$\bar w_H-1$ bounds the family's $\pi$-weighted value of compactness from below. Member $i$ has
$w_i\le(1+v_{N,i})/(1+m_D)$ if and only if $\rho_i\ge1+m_D$, that is, if it would train on more tokens at its own size
when it bears the data cost in proportion to its compute share, which a binding cap on the family's budget does not
imply. If every member would train on more tokens at its own size
($\rho_i\ge1+m_D$ for each $i$), each member's $w_i-1$ bounds its own $v_{N,i}$ from below (with $m_D=0$).
\end{enumerate}
\end{proposition}

\begin{proof}
(i) By Lemma~\ref{lem:geometry}(iii), $\ln w=\ln(a_1A/b_1B)-a_1n+b_1d$ in the generalized family. (ii) The
first-order condition in $d$ is $\sum_i\Lambda_i\varepsilon_{D,i}=\sum_iX_{T,i}+q_DD$. Divide it by $\sum_jX_{T,j}$.

(iii) The condition in $n_i$ is $\Lambda_i\varepsilon_{N,i}=X_{T,i}(1+m_{N,i})+\xi_i$, with $\xi_i=0$ at an interior
optimum, so $\rho_i=(1+v_{N,i})/w_i$; substitute into (ii) and multiply by $\bar w_H$. For the shares,
$1-\sum_ih_i/w_i=\sum_ih_i(1-1/w_i)$. With $m_D=0$, the $\pi$-weighted mean of $v_{N,i}$ is $\bar w_H-1$, which equals
the $\pi$-weighted mean of $w_i-1$ because $\sum_i\pi_iw_i=\bar w_H$; without binding size constraints
$v_{N,i}=m_{N,i}=\lambda\varrho T_i/(3D)$, and multiplying by $3D$ gives the statement on $\hat T_i$. (iv) Substitute
$v_{N,i}=m_{N,i}+\xi_i/X_{T,i}$ in (iii), with the sign of $\xi_i$. (v) Proposition~\ref{prop:A-wedge}(i) member by
member, and (i).

(vi) With the cap binding, the condition in $d$ becomes $\sum_ih_i\rho_i=1+m_D+\mu$; substituting
$\rho_i=(1+v_{N,i})/w_i$ from (iii) and multiplying by $\bar w_H$ gives the equality. Since $\rho_i=(1+v_{N,i})/w_i$,
$\rho_i\ge1+m_D$ if and only if $w_i\le(1+v_{N,i})/(1+m_D)$. A member that bears the share $h_i$ of the data cost gains
$\Lambda_i\varepsilon_{D,i}=\rho_iX_{T,i}$ and pays $X_{T,i}+h_iq_DD=X_{T,i}(1+m_D)$ per log point of tokens at its own
size, so $\rho_i\ge1+m_D$ says that it would train on more tokens.
\end{proof}

\begin{remark}
In two brute-force examples with three members on a common budget, the member inversions $\hat T_i/T_i$ are 0.86, 0.91
and 1.06 in one and 6.9, 2.3 and 0.65 in the other, while the $\pi$-weighted average is exact in both: member wedges can
be off by an order of magnitude. With the budget of the second example capped at 30, 50 and 80 percent of the family's
optimum, the equality in (vi) holds to numerical precision; one of the three members, the smallest, violates the member
bound at 30 and 50 percent, and two do at 80 percent, exactly those with $\rho_i<1$. With a binding size cap on one member
and a floor on another, the equalities of (iii) and (vi) in virtual values hold to numerical precision, with and without a
data cost and for the reference technology with $\kappa$ free. For the objective's own $m_{N,i}$, a size cap and a data
cap point in opposite directions: a capped member with $\rho_i\ge1$ has $w_i-1\le v_{N,i}$ but can have $w_i-1>m_{N,i}$
(\texttt{code/analysis/rb5\_theory/verify\_virtual\_value.py}).
\end{remark}

\subsection{Model-Free Identification of the Elasticity of Substitution and the Wedge}\label{app:modelfree}

\begin{proposition}[Model-free $\sigma^*$ and $w$; Proposition~\ref{prop:modelfree}]\label{prop:A-modelfree}
Let $y=F(n,d)$ be $C^2$ with $F_n,F_d>0$, where $y$ is any strictly increasing transformation of $-L$. Write
$x\equiv\ln M=d-n$ and $c\equiv\ln C=\ln6+n+d$, let $Y(x,c)\equiv F\big((c-\ln6-x)/2,\,(c-\ln6+x)/2\big)$, let $L$ also
denote loss as a function of $(x,c)$, and let $L^*(c)\equiv\min_xL(x,c)$ be the loss--compute frontier with argmin
$x^*(c)=\ln M^*(C)$.
\begin{enumerate}
\item[(i)] (Wedge.) At any point, $F_n=Y_c-Y_x$, $F_d=Y_c+Y_x$ and
\[
w=\frac{Y_c-Y_x}{Y_c+Y_x}=\frac{L_c-L_x}{L_c+L_x}.
\]
$w$ is invariant to strictly monotone transformations of output, so it requires neither $E$ nor $\kappa$ nor a
functional form: only the slopes of any output measure along the isocost and along the ray of constant $M$.
\item[(ii)] (Elasticity of substitution at the compute optimum.) At $x=x^*(c)$, where $L_x=0$,
\begin{equation}\label{eq:A-mf}
\frac1{\sigma^*}-1=\frac{2(-Y_{xx})}{Y_c}=\frac{2L_{xx}}{|dL^*/dc|}=\frac{L_{nn}|_C}{2\,|dL^*/dc|},
\end{equation}
where $L_{nn}|_C$ is the second derivative of the IsoFLOP profile with respect to $\ln N$ at its minimum and
$dL^*/dc=L_c$ by the envelope theorem. The ratio is invariant to strictly monotone transformations of output (with a nonzero
derivative) at the argmin, but not elsewhere: the curvature of the IsoFLOP profile divided by twice the slope of the loss--compute frontier,
in any common output units, identifies $\sigma^*$.
\item[(iii)] (Slope of the log wedge.) At $x^*(c)$, $\partial\ln w/\partial\ln M|_C=1/\sigma^*-1$. In the generalized
family $\ln w$ is exactly linear in $\ln(M/M^*(C))$ with this slope at every compute (Lemma~\ref{lem:geometry}(iii)); in
general the slope varies with $M$, which gives a test of the family.
\item[(iv)] (What $\kappa=1$ imposes.) In the generalized family, with $y=-\ln(L-E)$, $dy^*/dc=\gamma$ and
$-y_{nn}|_C=S\gamma$ at the argmin, so \eqref{eq:A-mf} returns $S/2$ for every $\kappa$ and $E$. The restriction $\kappa=1$
imposes $S=\gamma/[a(1-a)]$: the curvature of every IsoFLOP profile in log reducible loss is pinned by the slopes of the
frontier and of the path. With $\kappa$ free the curvature is a separate parameter, which \eqref{eq:A-mf} measures.
\item[(v)] (Local $\sigma$ anywhere.) $F_{nn}=Y_{cc}-2Y_{cx}+Y_{xx}$, $F_{dd}=Y_{cc}+2Y_{cx}+Y_{xx}$ and
$F_{nd}=Y_{cc}-Y_{xx}$, so $\sigma=\mathcal P/(\mathcal P+\mathcal Q)$ (Lemma~\ref{lem:soc}) is identified at any point from the second-order
expansion of any output in $(x,c)$.
\item[(vi)] (Designs.) $\sigma^*$ is identified by IsoFLOP profiles: the curvature of a profile at its minimum, and the
frontier slope from the minima at two or more budgets (or locally). On-path designs identify the frontier slope but not
the curvature. The wedge at a point is identified by local variation in two non-collinear directions around it, and only
inside the design's support; its precision falls as $\varepsilon_D\to0$, since $w\to\infty$ when $L_c+L_x\to0$.
\end{enumerate}
\end{proposition}

\begin{proof}
(i) $\partial/\partial n=\partial_c-\partial_x$ and $\partial/\partial d=\partial_c+\partial_x$. With $L=h(Y)$ and $h'<0$,
$L_x=h'Y_x$ and $L_c=h'Y_c$, and $h'$ cancels. (ii) At $F_n=F_d\equiv f$, $\mathcal P=2f^3$ and
$\mathcal Q=-f^2(F_{nn}-2F_{nd}+F_{dd})=-4f^2Y_{xx}$, since $Y_{xx}=(F_{nn}-2F_{nd}+F_{dd})/4$. By Lemma~\ref{lem:soc},
$1/\sigma-1=\mathcal Q/\mathcal P=-2Y_{xx}/f$ with $f=Y_c$. For $L=h(Y)$,
$L_{xx}=h'Y_{xx}$ at $Y_x=0$ and $dL^*/dc=h'Y_c$; at fixed $c$, $dn=-dx/2$, so $L_{nn}|_C=4L_{xx}$. For a monotone $g$,
$(g\circ L)_{xx}=g''L_x^2+g'L_{xx}=g'L_{xx}$ only where $L_x=0$, and the ratio is preserved where $g'\neq0$. (iii) $\partial_x\ln w=(Y_{cx}-Y_{xx})/(Y_c-Y_x)-(Y_{cx}+Y_{xx})/(Y_c+Y_x)$,
which is $-2Y_{xx}/Y_c$ at $Y_x=0$. (iv) Along the path $y^*=\gamma(c-\ln6)-\ln K$, so $dy^*/dc=\gamma$; by (ii),
$-y_{nn}|_C=-4Y_{xx}=S\gamma$; with $\kappa=1$, $\gamma=a(1-a)S$. (v) Second-order chain rule. (vi) Part (ii) uses only
derivatives of a profile at its minimum and of the frontier; part (i) uses two directional derivatives at the point.
\end{proof}

\begin{corollary}[Known costs that are not multiplicative; Section~\ref{sec:ident}]\label{cor:A-general}
Let $y=F(n,d)$ be $C^2$ with $F_n,F_d>0$, where $y$ is any strictly increasing transformation of $-L$, and let the cost
$C(N,D)$ be any known twice-differentiable function, with $\mathcal C(n,d)\equiv\ln C(e^n,e^d)$ and log-cost elasticities
$g_n\equiv\mathcal C_n>0$ and $g_d\equiv\mathcal C_d>0$. At a point that maximizes output on its isocost (equivalently,
minimizes cost at its output), let $\sigma^*$ be the technology's Hicks elasticity of substitution there, let
$r\equiv g_n/g_d$, let $\sigma_C$ be the Hicks elasticity of substitution of the
isocost itself (of $C$ read as a production function), let $y_{nn}|_C$ be the second derivative of output along the
isocost with respect to $n$, and let $y^*(c)$ be the output--cost frontier. Then
\begin{equation}\label{eq:A-general}
\frac1{\sigma^*}-\frac1{\sigma_C}=\frac{\lvert y_{nn}|_C\rvert}{r(1+r)\,g_d\,dy^*/dc}.
\end{equation}
\begin{enumerate}
\item[(i)] (Multiplicative cost.) With $C=6ND$, $\sigma_C=1$ and $r=g_d=1$, and \eqref{eq:A-general} is \eqref{eq:A-mf}.
\item[(ii)] (Linear cost.) With $C=q_NN+q_DD$, $\sigma_C=\infty$ and $(g_n,g_d)$ are the cost shares $(s_N,s_D)$, so
$1/\sigma^*=\lvert y_{nn}|_C\rvert\,s_D/(s_N\,dy^*/dc)$.
\item[(iii)] (Interior optima.) The second-order condition for a cost minimum holds strictly if and only if
$1/\sigma^*>1/\sigma_C$; for positive elasticities, $\sigma^*<\sigma_C$. This generalizes Lemma~\ref{lemma:sigma}
($\sigma_C=1$) and the condition $\sigma^*>0$ under a linear cost ($\sigma_C=\infty$).
\item[(iv)] (Invariance.) The right side of \eqref{eq:A-general} is invariant to strictly monotone transformations of
output with a nonzero derivative at the point, because the first derivative of output along the isocost vanishes there;
it may be computed from loss, with $L_{nn}|_C$ and $\lvert dL^*/dc\rvert$ in place of $\lvert y_{nn}|_C\rvert$ and
$dy^*/dc$.
\end{enumerate}
\end{corollary}

\begin{proof}
Parameterize the isocost as $d=\phi(n)$. Differentiating $\mathcal C(n,\phi(n))=c$ twice gives $\phi'=-r$ and
$\phi''=\mathcal Q_C/g_d^3$, with $\mathcal Q_C\equiv-(\mathcal C_{nn}g_d^2-2\mathcal C_{nd}g_ng_d+\mathcal C_{dd}g_n^2)$; by
Lemma~\ref{lem:soc}(i) applied to $\mathcal C$, $\mathcal Q_C=\mathcal P_C(1/\sigma_C-1)$ with
$\mathcal P_C\equiv g_ng_d(g_n+g_d)=r(1+r)g_d^3$. Output along the isocost, $Y(n)\equiv F(n,\phi(n))$, has
$Y'=F_n-rF_d=0$ at the point and $Y''=F_{nn}-2rF_{nd}+r^2F_{dd}+F_d\phi''$. With $F_n=rF_d$,
$F_{nn}-2rF_{nd}+r^2F_{dd}=-\mathcal Q/F_d^2$, and Lemma~\ref{lem:soc}(i) gives $\mathcal Q=\mathcal P(1/\sigma^*-1)$ with
$\mathcal P=r(1+r)F_d^3$. Hence $y_{nn}|_C=Y''=-r(1+r)F_d\,(1/\sigma^*-1/\sigma_C)$. By the envelope theorem $dy^*/dc$ is
the Lagrange multiplier, $F_d/g_d$, and dividing gives \eqref{eq:A-general}. (i) $\mathcal C$ is linear, so
$\mathcal Q_C=0$. (ii) With $\mathcal C=\ln(q_Ne^n+q_De^d)$, $g_n=s_N$, $g_d=s_D$ and
$\mathcal C_{nn}=\mathcal C_{dd}=-\mathcal C_{nd}=s_Ns_D$, so $\mathcal Q_C=-s_Ns_D=-\mathcal P_C$ and $1/\sigma_C=0$; and
$r(1+r)g_d=s_N/s_D$. (iii) $Y''<0$ if and only if $1/\sigma^*>1/\sigma_C$, since $r,F_d>0$. (iv) For $h$ strictly
monotone, $(h\circ Y)''=h''Y'^2+h'Y''=h'Y''$ at $Y'=0$ and $d\,h(y^*)/dc=h'\,dy^*/dc$. The identity was checked
symbolically for generic derivatives of $F$ and $\mathcal C$, and numerically, to at least five digits, for CES
technologies with $\sigma^*=0.40$, 0.67 and 0.83 under $C=ND$ and with $\sigma^*=0.40$, 0.67, 0.83 and 1.43 under a linear
cost, for the reference technology under multiplicative, linear and size-dependent-price costs, for a CES technology
under a CES cost, and for a non-separable technology with a data cost
(\texttt{code/analysis/rb5\_theory/verify\_general\_isocost.py}).
\end{proof}

\begin{remark}[Finite grids]\label{rem:A-grid}
IsoFLOP profiles are sampled on finite grids and fitted by low-order polynomials \citep{czech2026problems}. If the profile
$f(t)$, with $t$ the distance from its minimum in $\ln N$, is sampled at points symmetric about the minimum and fitted by
least squares on $(1,t,t^2)$, the quadratic coefficient is $f''/2+(f''''/24)\kappa_4+O(h^4)$, with
$\kappa_4\equiv(m_6-m_2m_4)/(m_4-m_2^2)$, $m_k$ the $k$-th moment of the grid and $h$ its half-width, and the fitted minimum
is displaced by $-(f'''/6)(m_4/m_2)/f''+O(h^3)$. For the \citet{besiroglu2024chinchilla} technology at $10^{21}$ FLOP and nine
equally spaced sizes, a grid spanning a factor of 9 (100) in $N$ overstates the curvature by 1.3 (5.9) percent, so the
model-free estimate is $\hat\sigma^*=0.734$ (0.726) against 0.737; a quartic fit removes the leading term. Quadratic fits
on grids narrower than a factor of about 3 in $N$ keep the bias below 0.001. These are noise-free calculations of bias
only: a narrower grid or a quartic fit raises the variance of the estimated curvature.
\end{remark}

\subsection{Partial Identification beyond the Design}\label{app:pi}

\begin{proposition}[Partial identification of $M^*(C)$ and of the wedge; Proposition~\ref{prop:modelfree}, beyond the design]\label{prop:A-pi}
Let $x^*(c)\equiv\ln M^*(C)$ be the unique argmin of the IsoFLOP profile at compute $C$, absolutely continuous in $c$, with
path slope $e_p(c)\equiv dx^*/dc=1-2a(c)$, where $a(c)\equiv d\ln N^*/d\ln C$. Suppose $x^*(c_J)$ is identified at the
largest design budget $c_J$, and consider a model with compute $c>c_J$ and log ratio $x=\ln M$.
\begin{enumerate}
\item[(i)] (Normal inputs.) If $N^*(C)$ and $D^*(C)$ are nondecreasing, then $e_p\in[-1,1]$ and
$x^*(c)\in[x^*(c_J)-(c-c_J),\ x^*(c_J)+(c-c_J)]$.
\item[(ii)] (Bounded path slope.) If $e_p(c')\in[e_L,e_U]$ for $c'\in[c_J,c]$, the sharp identified set is
\[
\mathcal X^*(c)=\big[x^*(c_J)+e_L(c-c_J),\ x^*(c_J)+e_U(c-c_J)\big],
\]
whose width grows linearly in the log extrapolation distance; monotonicity of $M^*$ is the case $e_L\ge0$.
When the true technology is in the generalized family, every point of $\mathcal X^*(c)$ is generated by a technology with
positive marginal products and single-peaked IsoFLOP profiles that coincides with the true one for $c\le c_J$. If
confidence intervals $\mathrm{CI}_j$ for $x^*(c_j)$ at budgets $c_1<\dots<c_J$ hold jointly with probability $1-p$ and the
slope bounds hold from $c_1$ on, then
$\bigcap_j[\inf\mathrm{CI}_j+e_L(c-c_j),\ \sup\mathrm{CI}_j+e_U(c-c_j)]$ covers $x^*(c)$ with probability at least $1-p$.
\item[(iii)] (Sign of the wedge.) If each IsoFLOP profile is single-peaked in $x$, that is, $Y_x>0$ for $x<x^*(c)$ and
$Y_x<0$ for $x>x^*(c)$, which holds whenever $0<\sigma<1$ throughout (Lemma~\ref{lem:soc}), then
$\operatorname{sign}(w-1)=\operatorname{sign}\big(x-x^*(c)\big)$ at every $(c,x)$, for any technology. Hence $w>1$ is
identified if $x>\sup\mathcal X^*(c)$ and $w<1$ if $x<\inf\mathcal X^*(c)$; when $\mathcal X^*(c)$ is sharp, as in (ii)
for a technology in the generalized family, these conditions are also necessary, and the sign is not identified
otherwise. Neither $\sigma^*$ nor the functional form is needed.
\item[(iv)] (Magnitude.) Suppose $\partial\ln w/\partial x|_c\in[\theta_L,\theta_U]$ with $\theta_L>0$ between $x^*(c)$ and
$x$; in the generalized family $\theta_L=\theta_U=1/\sigma^*-1$, so a range for $\sigma^*$ gives $[\theta_L,\theta_U]$.
With $\mathcal X^*(c)=[\ell,u]$: $\ln w\in[\theta_L(x-u),\,\theta_U(x-\ell)]$ if $x>u$;
$\ln w\in[\theta_U(x-u),\,\theta_L(x-\ell)]$ if $x<\ell$; and $\ln w\in[\theta_U(x-u),\,\theta_U(x-\ell)]$ if
$\ell\le x\le u$.
\item[(v)] (Compactness share and $T$.) Since $s=1-1/w$ is increasing in $w$, bounds $[w_L,w_U]$ give
$s\in[1-1/w_L,\,1-1/w_U]$. Under Proposition~\ref{prop:A-wedge}(iii)(a) with $\phi=1$ and $\eta=\delta=0$,
$\lambda\varrho T/(3D)\in[w_L-1,\,w_U-1]$.
\item[(vi)] (A bound inside the support.) Let $F_{nn},F_{dd}\le0$ and $F_{nd}\ge0$, which hold throughout the generalized
family. Then $\ln w$ is nonincreasing in $n$ and nondecreasing in $d$, so for any support point with $n_b\ge n_0$ and
$d_b\le d_0$, $w(n_0,d_0)\ge w(n_b,d_b)$, where $w(n_b,d_b)$ is identified without extrapolation. No finite upper bound
follows from these shape restrictions, and without $F_{nd}\ge0$ the bound can fail.
\item[(vii)] (Interval anchor, joint path bands and tilt allowance.) Let the anchored technology have $x^*(c_J)$ in
$[\underline x_J,\overline x_J]$ and path slopes in $[e_L,e_U]$ on $[c_J,c]$, and let the developer's compute-optimal log
ratio lie within $\ln\tau$ of the anchored technology's at every compute, $|x^*_{\mathrm{dev}}(c')-x^*(c')|\le\ln\tau$
with $\tau\ge1$. Then
\[
x^*_{\mathrm{dev}}(c)\in\mathcal X^*_\tau(c)\equiv\big[\underline x_J+e_L(c-c_J)-\ln\tau,\ \
\overline x_J+e_U(c-c_J)+\ln\tau\big],
\]
and, if the developer's IsoFLOP profiles are single-peaked, (iii) and (iv) hold for the wedge in its first-order
conditions with $\mathcal X^*_\tau(c)$ in place of $\mathcal X^*(c)$ and $x^*_{\mathrm{dev}}$ in place of $x^*$. If the
anchor is a point ($\underline x_J=\overline x_J$) and the anchored technology is in the generalized family,
$\mathcal X^*_\tau(c)$ is sharp: every point of it is $x^*_{\mathrm{dev}}(c)$ for a developer technology with positive
marginal products and single-peaked IsoFLOP profiles that lies within $\ln\tau$, at every compute, of a technology that
coincides with the anchored one for $c'\le c_J$ and has path slopes in $[e_L,e_U]$. With sampling error, let
$\mathcal B(c')=[b_L(c'),b_U(c')]$ be a simultaneous confidence band for the anchored path over $[c_J,\bar c]$ at level
$1-p$, $\Pr\{x^*(c')\in\mathcal B(c')\ \text{for all}\ c'\in[c_J,\bar c]\}\ge1-p$, for instance, for a path linear in
$c$, one built from a joint confidence set for its level and slope. Then $[b_L(c)-\ln\tau,\,b_U(c)+\ln\tau]$ covers $x^*_{\mathrm{dev}}(c)$ at
every $c\in[c_J,\bar c]$ simultaneously with probability at least $1-p$.
\item[(viii)] (Several anchored paths.) If the developer's compute-optimal log ratio lies within $\ln\tau$ of at least
one of finitely many anchored paths $j$, each with its own anchor budget and its own set $\mathcal X^*_j(c)$ from (vii)
with $\tau=1$, it lies in $\bigcup_j\mathcal X^*_{j,\tau}(c)$, and (iii) and (iv) hold with the union, with $\ell$ and $u$
its infimum and supremum. If each $\mathcal X^*_j$ is a simultaneous confidence band for path $j$ over its range of compute
up to $\bar c$ at level $1-p_j$, the widened union covers $x^*_{\mathrm{dev}}(c)$ with probability at least
$1-\max_jp_j$, because it covers whenever the band of the path that the developer's technology lies within $\tau$ of
covers.
\end{enumerate}
\end{proposition}

\begin{proof}
(i) $d\ln N^*+d\ln D^*=d\ln C$ with both terms nonnegative gives $a\in[0,1]$; integrate $e_p=1-2a$. (ii) Integrate $e_p$
over $[c_J,c]$; the linear continuations attain every point. For the technology: in the generalized family let
$\tilde L(n,d)=E+[Ae^{-a_1(n+\psi_N(c))}+Be^{-b_1(d+\psi_D(c))}]^\kappa$ with $c=\ln6+n+d$ and nondecreasing
$\psi_N,\psi_D$ that vanish for $c\le c_J$. At fixed $c$ the profile is a translate of a profile of the original
technology, hence single-peaked, with argmin $x^*(c)+2(1-a)\psi_N(c)-2a\psi_D(c)$ (Lemma~\ref{lem:alloc}(iii));
$\tilde F_n=F_n(1+\psi_N')+F_d\psi_D'>0$ and $\tilde F_d=F_n\psi_N'+F_d(1+\psi_D')>0$. Choosing
$\psi_N'=\Delta'_+/[2(1-a)]$ and $\psi_D'=\Delta'_-/(2a)$, with $\Delta'_\pm$ the positive and negative parts of $\Delta'$,
generates any continuation $x^*(c)+\Delta(c)$ with $\Delta(c_J)=0$. On the event that every $\mathrm{CI}_j$ covers, every
cone contains $x^*(c)$. (iii) By Proposition~\ref{prop:A-modelfree}(i), $w-1=-2Y_x/(Y_c+Y_x)=-2Y_x/F_d$ with $F_d>0$, and
single-peakedness gives the sign of $Y_x$. If $0<\sigma<1$ throughout, every critical point of a profile is a strict
local maximum of $y$ (Lemma~\ref{lem:soc}), so there is at most one. Sharpness makes every $x$ inside the set consistent
with both signs. (iv) $\ln w(x)=\int_{x^*}^x\partial_x\ln w$, with $\ln w(x^*)=0$ and $x^*\in[\ell,u]$. (v) Monotonicity.
(vi) $\partial\ln w/\partial n=F_{nn}/F_n-F_{nd}/F_d\le0$ and $\partial\ln w/\partial d=F_{nd}/F_n-F_{dd}/F_d\ge0$; move
from $(n_b,d_b)$ to $(n_b,d_0)$ and then to $(n_0,d_0)$. A saturating continuation of the data term beyond the support
satisfies the shape restrictions and sends $w\to\infty$; the concave quadratic
$F=0.05n+0.3d-0.025n^2-0.005d^2-0.02nd$ has $F_{nd}<0$ and $w(0,0.5)<w(0,0)$.

(vii) Integrating $e_p$ from a point of $[\underline x_J,\overline x_J]$ puts the anchored $x^*(c)$ in
$[\underline x_J+e_L(c-c_J),\,\overline x_J+e_U(c-c_J)]$, and the allowance widens it by $\ln\tau$ on each side; (iii)
and (iv) use only the developer's own technology. For sharpness, write a point of $\mathcal X^*_\tau(c)$ as $y_0+t$ with
$y_0$ in the set of (ii) and $|t|\le\ln\tau$. By (ii), $y_0$ is the compute-optimal log ratio at $c$ of a technology
$\tilde F$ with positive marginal products and single-peaked profiles that coincides with the anchored one for $c'\le c_J$
and has path slopes in $[e_L,e_U]$. The constant factor-augmenting tilt $\psi_N=-\psi_D=t/2$, whose bias index is
$\chi=-(a_1+b_1)t/2$, maps $(n,d)$ into $(n+t/2,d-t/2)$ and leaves compute unchanged, so at every compute the tilted
technology's IsoFLOP profile is that of $\tilde F$ translated by $t$ in $x$: its argmin moves by $t=-2\chi/(a_1+b_1)$ at
every compute, as in Lemma~\ref{lem:alloc}(iii), its profiles stay single-peaked, and its marginal products are those of
$\tilde F$ at shifted inputs, hence positive. It lies within $\ln\tau$ of $\tilde F$ at every compute and has
$x^*_{\mathrm{dev}}(c)=y_0+t$. On the event that $\mathcal B$ covers the anchored path on $[c_J,\bar c]$, the widened band
covers $x^*_{\mathrm{dev}}$ there. (viii) Let $j^*$ be a path that the developer's technology lies within $\tau$ of. By
(vii), $x^*_{\mathrm{dev}}(c)\in\mathcal X^*_{j^*,\tau}(c)$, which is contained in the union, and the event that band $j^*$
covers has probability at least $1-p_{j^*}\ge1-\max_jp_j$. The tilt construction, the containment and sign statements
(20,000 random draws) and the band coverage (a Monte Carlo with three anchored paths) were also checked numerically
(\texttt{code/analysis/rb5\_theory/verify\_prop\_pi.py}).
\end{proof}

\begin{remark}[Parametric projection and the direction of extrapolation errors]\label{rem:A-pi-param}
A confidence set for $(\ln(A/B),\alpha,\beta)$ projects into one for the wedge whose width grows linearly in the log
extrapolation distance, since $\Var(\ln\hat w)$ is exactly quadratic in $\ln M$ at fixed compute. For the Chinchilla
extraction a pairs bootstrap gives $\mathrm{sd}(\ln\hat w)=0.161$ at $M=1{,}875$ and $C=7.2\times10^{23}$, which is sampling
error within one sweep and one functional form; the union of the designs' path bands of
Proposition~\ref{prop:A-pi}(viii), which carries each path's slope and its sampling error, puts $M^*(10^{24})$ in a set
about 3.8 log points wide (Online Appendix~\ref{app:wedge-pi}), and the 32
technologies in the set put $M^*(10^{24})$ between 1.4 and 426, or 5.7 log points. Evidence that
fitted laws misstate losses at extreme ratios \citep{sardana2024beyond} does not sign the error in $\hat w$: revealed
preference recovers the technology the developer believed (Proposition~\ref{prop:A-wedge}(ii)), so if developers planned with
fitted laws, correcting toward the true technology moves $\hat w$ away from their first-order condition.
\end{remark}
